\SourceLeaf{376}%
% A544-P0376-D01
\Article{16}
Es ist folglich
\[
\text{Dec. }\frac mM-\text{Dec. }\frac Mm=
\]
\[
\begin{aligned}
\mathrm{I.}\;&-4\Sigma\varepsilon,\ \text{wo }y,Y\text{ ganz, }[x],[X]\text{ gerade}\\
\mathrm{II.}\;&+4\Sigma\varepsilon,\ Y\text{ ganz, }[X]\text{ gerade, }[x]\text{ gerade, }[y]\text{ ungerade}\\
\mathrm{III.}\;&-4\Sigma\varepsilon,\ x\text{ ganz, }[X]\text{ gerade, }[Y]\text{ ungerade,}\\
&\qquad [y]\text{ ungerade}\\
\mathrm{IV.}\;&+4\Sigma\varepsilon,\ Y\text{ ganz, }[X]\text{ gerade, und nicht zugleich }[x],[y]\text{ gerade}\\
\mathrm{V.}\;&-4\Sigma\varepsilon,\ x\text{ ganz, }[y]\text{ gerade und nicht zugleich }[X],[Y]\text{ gerade}\\
&+\psi
\end{aligned}
\]

% A544-P0376-D02
Hier ist $\psi$ in folgender Tabelle dargestellt
\begingroup\small\setlength{\arraycolsep}{1.8pt}
\[
\begin{array}{cccccc|rrrr||cccccc|rrrr}
a&b&A&B&\alpha&\beta&\multicolumn{4}{c||}{\psi}&a&b&A&B&\alpha&\beta&\multicolumn{4}{c}{\psi}\\
+&+&+&+&+&+&+4&0&0&-2&-&-&+&+&-&-&0&0&0&-2\\
&&+&+&+&-&0&0&0&-2&&&+&+&-&+&+4&0&0&-2\\
&&-&+&+&-&0&-4&0&-2&&&-&+&-&+&0&-4&0&-2\\
&&-&+&-&-&0&-4&0&-2&&&-&+&+&+&0&-4&0&-2\\
&&-&-&-&-&-4&-4&0&-2&&&-&-&+&+&0&-4&0&-2\\
&&-&-&-&+&0&-4&0&-2&&&-&-&+&-&-4&-4&0&-2\\
&&+&-&-&+&0&0&0&-2&&&+&-&+&-&0&0&0&-2\\
&&+&-&+&+&0&0&0&-2&&&+&-&-&-&0&0&0&-2\\[4pt]
-&+&+&+&-&+&+4&0&0&-2&+&-&+&+&+&-&0&0&0&-2\\
&&+&+&+&+&+4&0&0&-2&&&+&+&-&-&0&0&0&-2\\
&&-&+&+&+&0&-4&0&-2&&&-&+&-&-&0&-4&0&-2\\
&&-&+&+&-&0&-4&0&-2&&&-&+&-&+&0&-4&0&-2\\
&&-&-&+&-&-4&-4&0&-2&&&-&-&-&+&0&-4&0&-2\\
&&-&-&-&-&-4&-4&0&-2&&&-&-&+&+&0&-4&0&-2\\
&&+&-&-&-&0&0&0&-2&&&+&-&+&+&0&0&0&-2\\
&&+&-&-&+&0&0&0&-2&&&+&-&+&-&-4&-4&-4&-6\\
\end{array}
\]
\endgroup

% A544-P0376-D03
\[
\begin{array}{lll}
\text{Hier gelten die ersten beiden Columnen}&\text{für }\frac12(M-1)\text{ gerade}\\
\text{letzten beiden}&\text{für }\frac12(M-1)\text{ ungerade}\\
\text{die erste und dritte}&\text{für }\frac12(m-1)\text{ gerade}\\
\text{zweite und vierte}&\text{für }\frac12(m-1)\text{ ungerade}
\end{array}
\]

\SourceLeaf{377}%
% A544-P0377-D01
\Article{17}
Die $128$ Fälle, welche in obiger Tafel bei der Bestimmung von $\psi$ unterschieden sind, lassen sich viel kürzer auf folgende Weise umfassen:
\[
\psi=k+l
\]
$k=-4$, wenn zugleich $a,A,\alpha,b,B,\beta$ die Zeichen $+\ +\ +\ -\ -\ -$ haben, sonst immer
\[
k=0
\]

% A544-P0377-D02
\begingroup\small\setlength{\arraycolsep}{3pt}
\[
\begin{array}{ll|ll}
\frac12(M-1)\text{ gerade}&\frac12(m-1)\text{ gerade}&l=+4,&\text{wenn }AB\beta\text{ positiv}\\
&&-4,&\text{wenn }AB\beta\text{ negativ}\\
&&0,&\text{in allen übrigen Fällen}\\[3pt]
\frac12(M-1)\text{ gerade}&\frac12(m-1)\text{ ungerade}&l=-4,&\text{wenn }A\text{ negativ}\\
&&0,&\text{wenn }A\text{ positiv}\\[3pt]
\frac12(M-1)\text{ ungerade}&\frac12(m-1)\text{ gerade}&l=0&\\[3pt]
\frac12(M-1)\text{ ungerade}&\frac12(m-1)\text{ ungerade}&l=-2&
\end{array}
\]
\endgroup

% A544-P0377-D03
Zu versuchen ist noch, ob es vortheilhafter ist, $A$ und $a$ positiv, dagegen aber $m\equiv1$, $M\equiv1$ nur nach mod.\ $2$ (nicht nach Modulus $2+2i$) zu nehmen. Das Endresultat muss werden
\[
\text{Dec. }\frac mM-\text{Dec. }\frac Mm\equiv
\]
\[
\begin{array}{r|rrrr}
&\multicolumn{4}{c}{M\equiv}\\
m\equiv&1&1+2i&3&3+2i\\\hline
1&0&0&0&0\\
1+2i&0&2&2&0\\
3&0&2&0&2\\
3+2i&0&0&2&2
\end{array}
\]
Alles nach Mod.\ $4$.
\BookRule{.3\textwidth}

\SourceLeaf{378}%
% A544-P0378-D01
\par\bigskip\Needspace{9\baselineskip}
\pdfbookmark[1]{Theorie der biquadratischen Reste. VI.}{nachlass-biquadratic-VI}
\centerline{[VI.]}\medskip
\begin{center}\large\scshape Theorie der biquadratischen Reste.\end{center}
\BookRule{.18\textwidth}
\renewcommand{\Article}[1]{\par\medskip\Needspace{4\baselineskip}\pdfbookmark[2]{Art. #1}{nachlass-biquadratic-VI-art-#1}\centerline{\textbf{#1.}}\nobreak\smallskip}
\Article{1}
Eine Reihe ganzer complexer Zahlen $\varphi,\varphi',\varphi''$ u.\,s.\,w. sei so beschaffen, dass erstlich sie unter einander alle incongruent sind nach dem Modulus $\mu=\alpha+\beta i$, $\alpha$ und $\beta$ ganze reelle Zahlen bezeichnend, zweitens dass jede ganze complexe Zahl einer von jenen nach dem Modulus $\mu$ congruent ist. Unter dieser Voraussetzung heisst der Inbegriff der Zahlen $\varphi,\varphi',\varphi''$ u.\,s.\,w. ein vollständiges Restsystem für den Modulus $\mu$. Es ist bewiesen, dass die Anzahl der darin begriffenen Zahlen der Norm von $\mu$, d.\,i. der Zahl $\alpha\alpha+\beta\beta$ gleich ist, welche mit $\nu$ bezeichnet werden soll.

% A544-P0378-D02
\Article{2}
Unter den Zahlen $\varphi,\varphi',\varphi''$ u.\,s.\,w. ist Eine durch $\mu$ theilbare; wird dieselbe ausgeschlossen und der Inbegriff der übrigen mit $\chi$ bezeichnet, so bildet $\chi$ ein vollständiges System der durch den Modulus untheilbaren Reste, deren Anzahl $=\nu-1$. Beschränken wir die Untersuchung auf ungerade Modulen, so ist $\nu-1$ durch $4$ theilbar. Es werden dann ferner $f,if,-f,-if$ unter sich incongruent sein, folglich diejenigen Zahlen in $\chi$, welche resp.\ denen $if,-f,-if$ congruent sind, unter sich und von $f$ verschieden. (Associirte und zusammengesetzte Zahlen.)

% A544-P0378-D03
Hieraus ergiebt sich eine Zerlegung von $\chi$ in vier Gruppen oder partielle Systeme $\kappa,\kappa',\kappa'',\kappa'''$. Man setzt eine beliebige Zahl aus $\chi$, z.\,B.\ $\varphi$ in die Gruppe $\kappa$, und die drei den Zahlen $i\varphi,-\varphi,-i\varphi$ congruenten Glieder von $\chi$, der Reihe nach in die Gruppen $\kappa',\kappa'',\kappa'''$. Nachdem diese vier Zahlen aus $\chi$ gestrichen sind, setzt man eine beliebige der übrigen wieder in $\kappa$, und die drei auf ähnliche Art davon abhängigen in $\kappa',\kappa'',\kappa'''$. So fährt man fort, bis das ganze System $\chi$ vertheilt ist. Die Gruppen $\kappa,\kappa',\kappa'',\kappa'''$ sollen zusammengehörige Viertelsysteme heissen. Es ist klar, dass sie folgende Eigenschaften haben:

\SourceLeaf{379}%
% A544-P0379-D01
\begin{enumerate}
\renewcommand{\labelenumi}{\arabic{enumi})}
\item Jedes Viertelsystem besteht aus $\frac14(\nu-1)=\frac14(\alpha\alpha+\beta\beta-1)$ Zahlen.
\item Das Charakteristische eines Viertelsystems ist, dass keine der darin befindlichen Zahlen weder selbst, noch ihr Product in $i,-1$, oder $-i$, einer andern aus demselben Viertelsystem congruent ist, jede durch $\mu$ nicht theilbare Zahl aber, entweder selbst oder ihr Product durch $i,-1$, oder $-i$ sich darin findet, oder einer daraus congruent ist.
\item So wie aus der Multiplication der Zahlen in $\kappa$ mit $i,-1$ und $-i$, resp.\ die Zahlen in $\kappa',\kappa'',\kappa'''$ entstehen, oder ihnen congruente, so reproducirt die Multiplication der Zahlen in $\kappa'$, mit jenen Factoren, resp.\ die Zahlen in $\kappa'',\kappa''',\kappa$; die Multiplication der Zahlen $\kappa''$ reproducirt auf ähnliche Weise die Zahlen $\kappa''',\kappa,\kappa'$; endlich die Multiplication der Zahlen $\kappa'''$ reproducirt $\kappa,\kappa',\kappa''$.
\end{enumerate}

% A544-P0379-D02
Kürze halber kann diese gegenseitige Abhängigkeit der vier Viertelsysteme so ausgedrückt werden
\[
\begin{aligned}
\kappa'&\equiv i\kappa,\\
\kappa''&\equiv-\kappa\equiv i\kappa',\\
\kappa'''&\equiv i\kappa''\equiv-\kappa'\equiv-i\kappa.
\end{aligned}
\]

% A544-P0379-D03
\Article{3}
Wenn $m$ eine complexe ganze Zahl bedeutet, die mit $\mu$ keinen gemeinschaftlichen Divisor hat, und die sämmtlichen Zahlen eines Viertelsystems $\kappa$ mit $m$ multiplicirt werden, so bilden die Producte, oder beliebige ihnen congruente Zahlen ihrerseits auch ein Viertelsystem; und ebenso entstehen durch Multiplication der Zahlen der Systeme $\kappa',\kappa'',\kappa'''$ noch drei Viertelsysteme, die mit jenem zusammengehören werden. Der Beweis ist leicht zu führen. Diese vier neuen Systeme mögen mit $m\kappa,m\kappa',m\kappa'',m\kappa'''$ bezeichnet werden, gleichviel, ob die Producte selbst oder nur ihnen congruente Zahlen gewählt werden. Im letztern Fall kann dies so geschehen, dass man immer nur solche wählt, die sich in einem der Systeme $\kappa,\kappa',\kappa'',\kappa'''$ befinden. Auf diese Art ist also das System $\chi$, wenigstens allgemein zu reden, auf zwei verschiedene Arten in Viertelsysteme zerlegt. Nehmen wir an, dass $m\kappa$ gemeinschaftlich hat

% A544-P0379-D04
\[
\begin{array}{rl}
\text{mit }\kappa\dotsfive&\lambda\text{ Zahlen}\\
\kappa'\dotsfive&\lambda'\text{ Zahlen}\\
\kappa''\dotsfive&\lambda''\text{ Zahlen}\\
\kappa'''\dotsfive&\lambda'''\text{ Zahlen}
\end{array}
\]
so wird auch $\kappa'$ mit $m\kappa'$, $\kappa''$ mit $m\kappa''$, $\kappa'''$ mit $m\kappa'''$ gemein haben $\lambda$ Glieder; %
\SourceLeaf{380}%
% A544-P0380-D01
$\kappa''$ mit $m\kappa'$, $\kappa'''$ mit $m\kappa''$, $\kappa$ mit $m\kappa'''$, $\lambda'$ Glieder u.\,s.\,w. Es sei $\varepsilon$ der kleinste Rest von $\lambda'+2\lambda''+3\lambda'''$ nach dem Modulus $4$, oder $\varepsilon$ eine der vier Zahlen $0,1,2,3$, je nachdem $\lambda'+2\lambda''+3\lambda'''$ von der Form $4n,4n+1,4n+2,4n+3$ ist. Ich behaupte nun, dass $\varepsilon$ von der Anordnung des Viertelsystems $\kappa$ unabhängig ist.

% A544-P0380-D02
Um die Beweisführung zu erleichtern, bediene ich mich folgender Bezeichnung. $\Pi\psi$ soll die Zahl $0,1,2,3$ bezeichnen, je nachdem die durch $\mu$ nicht theilbare Zahl $\psi$ sich (selbst oder durch Congruenz Repräsentation) in der Gruppe $\kappa,\kappa',\kappa'',\kappa'''$ befindet. Von selbst hat man daher die Folge

% A544-P0380-D03
I.\quad $\Pi(i\psi)\equiv1+\Pi\psi\ (\text{mod. }4)$.

II.\quad Die Zahl $i^{-\Pi\psi}\pdot\psi$ findet sich, entweder selbst oder durch Congruenz Repräsentation in der Gruppe $\kappa$.

III.\quad $\Sigma\Pi m\varphi\equiv\varepsilon,\ (\text{mod. }4)$, wenn die Summation über alle in $\kappa$ befindliche Glieder $\varphi$ erstreckt wird.

% A544-P0380-D04
Es sei nun $k$ ein anderes Viertelsystem, bestehend aus $f,f',f''$ u.\,s.\,w., während $\kappa$ aus $\varphi,\varphi',\varphi''$ u.\,s.\,w. besteht. Ich setze voraus, was erlaubt ist, da die \emph{Ordnung} der Glieder in $\kappa$ willkürlich, dass $f$ mit $\varphi$ identisch oder zusammenhängend ist, $f'$ mit $\varphi'$, $f''$ mit $\varphi''$ u.\,s.\,w. Die mit $k$ zusammenhängenden Viertelsgruppen seien $k'(\equiv ik)$, $k''(\equiv-k)$, $k'''(\equiv-ik)$. Es habe ferner die Charakteristik $\mathrm{P}$ in Beziehung auf die Gruppen $k,k',k'',k'''$ dieselbe Bedeutung wie $\Pi$ in Beziehung auf $\kappa,\kappa',\kappa'',\kappa'''$, so dass $\mathrm{P}\psi=0,1,2,3$, je nachdem $\psi$ zu $k,k',k'',k'''$ gehört.

% A544-P0380-D05
Es wird demnach, wenn man von der Vertheilung der $\chi$ in die Viertelsysteme $k,k',k'',k'''$ anstatt von $\kappa,\kappa',\kappa'',\kappa'''$ ausgeht, an die Stelle von $\varepsilon$ treten der kleinste Rest von $\mathrm{P}mf+\mathrm{P}mf'+\mathrm{P}mf''+\mathrm{P}mf'''$ u.\,s.\,w. u.\,s.\,w. oder von $\Sigma\mathrm{P}mf$ nach dem Modulus $4$ und es handelt sich, zu beweisen, dass $\Sigma\mathrm{P}mf-\Sigma\Pi m\varphi$ durch $4$ theilbar ist.

% A544-P0380-D06
Wir schreiben diese Grösse so
\[
\begin{array}{l}
\phantom{+}\mathrm{P}mf+\mathrm{P}mf'+\mathrm{P}mf''+\mathrm{P}mf'''+\text{u.\,s.\,w.}\\
-\mathrm{P}m\varphi-\mathrm{P}m\varphi'-\mathrm{P}m\varphi''-\mathrm{P}m\varphi'''-\\
+\mathrm{P}m\varphi+\mathrm{P}m\varphi'+\mathrm{P}m\varphi''+\mathrm{P}m\varphi'''+\\
-\Pi m\varphi-\Pi m\varphi'-\Pi m\varphi''-\Pi m\varphi'''-
\end{array}
\]

% A544-P0380-D07
Da der Voraussetzung nach $f$ und $\varphi$ congruent sind oder zusammengehören, so gilt dasselbe auch von $mf$ und $m\psi$ und man hat

